\newpage

\section{\huge \textcolor{blue}{一阶方程换元法}}

\subsection{\LARGE 尺度变换}
在具体说明一阶方程换元法之前, 
先介绍尺度变换(即通过改变坐标来改变坐标轴的尺度)\\
令
\[ x_1 = \frac{x}{a}  , y_1 = \frac{y}{b} \]
作用:
\begin{enumerate}
    \item 改变单位
    \item 让变量无量纲化
    \item 减少或简化方程中的常数
\end{enumerate}

\subsection{\LARGE 伯努利方程}
下以伯努利方程为例:
{\Large
\begin{equation}
y' = p(x) y + q(x) y^n \hspace*{1em} 
(n \neq 0 \hspace*{0.5em}and\hspace*{0.5em} 1)
\end{equation}
}
同除$y_n$:
\[ \frac{y'}{y^n} = \frac{p(x)}{y^{n-1}} + q(x)  \]

记 $v = 1/y^{n-1}$, 得 $v' = (1-n) y' / y^n$
{\Large \[ \therefore \hspace*{1em} \frac{v'}{1-n} = p(x) v + q(x)  \] }
后化为标准线性方程即可解.

\subsection{\LARGE 一阶齐次ODE}
关于齐次的关键就是变形找到(或通过换元创造)如下形似:
{\Large
\begin{equation}
y' = F(y/x) 
\end{equation}
}
eg: $y' = \frac{x^2 y}{x^3 + y^3} = \frac{\frac{y}{x}}{1 + (\frac{y}{x})^3} $\\
通过观察发现齐次方程的不变性质:\\
\hspace*{2em}缩放变换保持一致(即同比例改变各轴尺度)
\[ x \rightarrow a x_1 , y \rightarrow a y_1  \]
\[ dx \rightarrow a dx_1 , dy \rightarrow a dy_1  \]

{\Large \[ \therefore \hspace*{1em} \frac{dy_1}{dx_1} = F(y/x) = \frac{dy}{dx}  \] }
下对方程$y' = F(y/x) $处理:\\
令 $z = y/x $, 而为了避免对除法求导导致计算量过大, 
故\\
变形为 $y = z x$, 有 $y' = z' x + z $
\[ \therefore \hspace*{1em} z' x + z = F(z)  \]
\[ \therefore \hspace*{1em} x \frac{dz}{dx} = F(z) - z \]
下给出一个例子:\\
已知: $y' = tan(\alpha + \pi /4)$ , $tan(\alpha) = \frac{y}{x}$, 求原曲线轨迹\\
解:
\[ y' = tan(\alpha + \pi /4 ) = \frac{tan(\alpha) + tan(\pi /4)}{1 - tan(\alpha) tan(\pi /4)} \]
\[ y' = \frac{\frac{y}{x} + 1}{1 - \frac{y}{x} }  \]
令$z = \frac{y}{x} $, 得 $y' = z' x + z$
\[ \therefore \hspace*{1em} z' x + z = \frac{z + 1}{1 - z}  \]
\[ x \frac{dz}{dx} = \frac{z + 1}{1 - z} - z = \frac{1 + z^2}{1-z}\]
\[ \frac{1 - z}{1 + z^2} dz = \frac{dx}{x} \]
对两侧同时积分:
\[ \therefore \hspace*{1em} arctan(z) - \frac{1}{2} ln(1 + z^2) = ln(x) + C  \]
\[ \therefore \hspace*{1em} arctan(\frac{y}{x}) = ln(\sqrt{x^2 + y^2} ) + C  \]
极坐标转化:\\
令$r = \sqrt{x^2 + y^2} $, $ \theta = arctan(\frac{y}{x}) $
\[ \theta = ln(r) + c\]
\[ \therefore \hspace*{1em} r = c_1 e^{\theta}  \]
